% ===================================================================
%  10. TAULA — Itsasoa. — Portua.
%  Traduction 2026 d'apres texto/io, controlee sur texto/fr, texto/en
%  et texto/es. Consigne : voir texto/eu/10-taula-01.tex.
%
%  Le francais decale sa numerotation d'un cran a partir de l'alinea 8.
%  C'est une coquille de l'imprimeur francais ; le basque suit l'ido.
%
%  LE TABLEAU DE LA MER, ET C'EST L'INVERSE EXACT DE CELUI DE LA FERME.
%  Le basque a bien ses propres mots de mer — « itsasoa »,
%  « ontzia », « arraunak », « arrantzalea », « portua » remonte au
%  latin comme partout — mais toute la marine militaire et industrielle
%  du XIXe siecle lui est arrivee toute faite, et il l'ecrit comme les
%  autres : « korazatua », « torpedoa », « fragata », « dikea »,
%  « helizea ». Le meme critere qu'en ukrainien : l'age du mot, non la
%  famille de la langue.
% ===================================================================

%%K t10-apar-1 apar
\VUsaut{4.0mm}
\VUtitre{15.0pt}{60}{10. TAULA}
\VUsaut{3.0mm}
\VUpk{11.4pt}{40}{Itsasoa. --- Portua.}
\VUsaut{3.0mm}

%%K t10-01-1 p
1. --- Udan, batzuek mendietan bakea eta freskotasuna bilatzen duten bitartean,
beste batzuek itsasertzera joatea nahiago dute. Hala egin genuen guk iaz,
\VUgras{ozeanoa} \textsuperscript{(1)} oso maite baitugu.

%%K t10-02-1 p
2. --- Niri dagokidanez, atsegin handia ematen dit ur-zabal neurrigabe horri
begiratzeak, \VUgras{zeruertzeraino} \textsuperscript{(2)} hedatua, eta ikusteak
nola abiatzen diren \VUgras{ontzi} erraldoiak \textsuperscript{(3)}
\VUgras{bidaia} luzera, Ameriketara doazen bidaiariekin.

%%K t10-03-1 p
3. --- Horregatik aitak, nire gustuak ezagutzen dituenak, beti hautatzen du
oporretarako gure \VUgras{gerra-portu} handietako baten ondoko itsasertz oso
lasai bat.

%%K t10-03-2 p
Gure azken bisitan \VUgras{eskuadra} etorri eta aingura bota zuen
\VUgras{olatu-hausle} \textsuperscript{(4)} erraldoiaren atzean, zeinak
\VUgras{gerra-kaia} \textsuperscript{(5)} ixten eta babesten baitu, eta aski
miretsi ahal izan nituen \VUgras{gerraontzi} desberdinak:
\VUgras{urpekuntzi} txikitik \textsuperscript{(10)}, zeina murgildu eta urpean
desagertzen baita, \VUgras{korazatu} erraldoiraino \textsuperscript{(9)}, zeinak
orain arte \VUgras{torpedo-ontziaren} \textsuperscript{(8)} erasoak baino ez
baititu beldur izan.

%%K t10-tit-2 sub
\VUsaut{3.0mm}
\VUpk{10.2pt}{40}{Ontzi txikiak.}
\VUsaut{3.0mm}

%%K t10-04-1 p
4. --- Azken honek \VUgras{trebeki} aurkitzen zekien jada burdinazko oskol lodiaren
gune ahula eta han \VUgras{torpedo} arriskutsua jartzen, zeinaren leherketak
ontzia hondoratzen baitu; baina hala ere urpekuntzia baino gutxiago beldur izan
behar zitzaion, kontrajarleek erraz ekiten baitiote atzetik.

%%K t10-05-1 p
5. --- \VUgras{Gurutzaontzi} korazatu azkarrak ere ikusi ahal izan nituen
\textsuperscript{(6)}, ia korazatu benetakoa bezain zaurgaitzak,
\VUgras{garraio-ontzi} batzuk \textsuperscript{(12)}, hiru edo lau
\VUgras{kanoiontzi} \textsuperscript{(7)} eta, azkenik, \VUgras{fragata}
\textsuperscript{(11)}. \VUgras{Ikuspegirik ederrena aingura altxatzen ari den
flota honena da.}

%%K t10-06-1 p
6. --- Zeinen bestelakoa den altzairuzko mole horien eta \VUgras{hiru mastadun}
dotoreen edo \VUgras{belaontzi} liraineen \textsuperscript{(13)} egitura, zeinak
merkataritza-flotan oraindik erabiltzen baitira eta zeinak bela guztiak zabalik
portuan sartzen ikusi bainituen \VUgras{kai-mutur} batetik
\textsuperscript{(14)} paseatzen. Egia da asko galtzen zutela \VUgras{haizea}
kontra zutenean. Bela guztiak jaso behar zituzten dikean sartzeko, eta
\VUgras{atoiontzi} \textsuperscript{(16)} bat behar zen haiek hartu eta kaira
eramateko, \VUgras{gabarra} arrunt batzuk bezala \textsuperscript{(17)}.

%%K t10-tit-3 sub
\VUsaut{3.0mm}
\VUpk{10.2pt}{40}{Merkataritza-portua.}
\VUsaut{3.0mm}

%%K t10-07-1 p
7. --- Aipatzen ari naizen portua horregatik da interesgarria: ez du gerra-kaia
bakarrik, lan erraldoiez artifizialki sortua, baizik eta \VUgras{merkataritza-portu}
zoragarri bat ere bai ibai handi baten \VUgras{bokalean}.

%%K t10-08-1 p
8. --- Bi dikeak banatzen dituen lur-zerrenda gotortuta eta babestuta dago
itsasora aurreratutako \VUgras{bateria} eta \VUgras{baluarte} lerro batez.
\VUgras{Harresietan} \textsuperscript{(15)} paseatzeko baimen berezia behar da.
Erraz lortu nuen nire osabaren bidez, zeina ingeniari baita
\VUgras{ontziolan} \textsuperscript{(18)}, eta ikustera joan nintzen nola
eraikitzen dituzten ontziak teilape erraldoien azpian.

%%K t10-09-1 p
9. --- Korazatu baten uretaratzea ere ikusi nuen, eta gogoan dut jendetza osoa
hartu zuen hunkidura, mole erraldoiak, bere buruari utzita, bere ohetik irrist
egin eta ibaiaren uretan agertu zenean.

%%K t10-10-1 p
10. --- Ontziolara iristeko, \VUgras{arma-plaza} \textsuperscript{(19)}
zeharkatzen da, non goarnizioko tropak egunero trebatzen baitira, eta burdinazko
\VUgras{zubitxotik} \textsuperscript{(20)} igarotzen, zeinak plaza
\VUgras{arsenalarekin} \textsuperscript{(21)} lotzen baitu.

%%K t10-tit-4 sub
\VUsaut{3.0mm}
\VUpk{10.2pt}{40}{Arsenala eta dikeak.}
\VUsaut{3.0mm}

%%K t10-11-1 p
11. --- Osabarekin ikusi nuen establezimendu handi hori, \VUgras{kanoiz}
\textsuperscript{(22)}, \VUgras{bala biribilez} \textsuperscript{(23)} eta
\VUgras{jaurtigaiz} betea. \VUgras{Galdaretegia} ere ikusi nuen
\textsuperscript{(24)}, non ontzientzako \VUgras{galdarak}
\textsuperscript{(25)} egiten baitituzte.

%%K t10-12-1 p
12. --- Ondo-ondoan \VUgras{dikeak} daude \textsuperscript{(26)}, dike lehorrak,
eta oso aipagarriak diren \VUgras{dike flotagarriak} \textsuperscript{(27)},
zeinetan gertutik ikusi ahal izan bainituen konpontzen ari zen ontzi batean
\VUgras{helizea} \textsuperscript{(28)}, \VUgras{gila} \textsuperscript{(29)} eta
ontziaren \VUgras{lema} \textsuperscript{(30)}.

%%K t10-13-1 p
13. --- Merkataritza-portua gerrakoa bezain aztergarria da, eta ordu asko igaro
nituen kaietan noraezean, non ibaian gora doazen \VUgras{gurpil-ontziak}
\textsuperscript{(31)} baitaude, eta \VUgras{garabi} bat
\textsuperscript{(32)} nola aritzen den begira, zama izugarriak lumatxoak balira
bezala jasoz.

%%K t10-tit-5 sub
\VUsaut{4.0mm}
\VUpk{10.2pt}{40}{Praktikoa.}
\VUsaut{3.0mm}

%%K t10-14-1 p
14. --- \VUgras{Ontzi transatlantikoak} \textsuperscript{(34)} miresten nituen,
kaialdea uzten \VUgras{bandera} \textsuperscript{(35)} astinduekin, \VUgras{praktiko}
trebe batek \textsuperscript{(36)} gidatuta, zeinak zoragarri jakiten baitzuen
\VUgras{ubidean} eusten, \VUgras{buiez} \textsuperscript{(40)} eta
\VUgras{seinalez} markatua, \VUgras{txalupa handiak} \textsuperscript{(41)}
saihestuz, zeinak hain zailak baitira gidatzen bela laukidun bakarraren azpian.
\VUgras{Brankan} \textsuperscript{(37)} bigarren mailako \VUgras{kamarotak}
\textsuperscript{(39)} bereizten nituen, eta \VUgras{txopan}
\textsuperscript{(38)} lehen mailako bidaiarien saloiak.

%%K t10-15-1 p
15. --- Batzuetan \VUgras{zama-ontziaren} \textsuperscript{(42)} zamaketari ere
begiratzen nion, zeinetara \VUgras{biltegitik} \textsuperscript{(43)} eta
\VUgras{portuko geltokitik} \textsuperscript{(44)} salgaiak isurtzen baitzituzten,
\VUgras{kaiko trenbideko} \textsuperscript{(45)} tren ugarik ekarriak.

%%K t10-tit-6 sub
\VUsaut{3.0mm}
\VUpk{10.2pt}{40}{Txalupan ibiltzea.}
\VUsaut{3.0mm}

%%K t10-16-1 p
16. --- Sarritan ibiltzen nintzen txalupan \VUgras{dike itxian}
\textsuperscript{(53)}, non \VUgras{txalupa} txikiak \textsuperscript{(46)}
babesten baitira, eta poza zen niretzat \VUgras{arraun} bat
\textsuperscript{(47)} eustea. \VUgras{Bizkarrera} \textsuperscript{(48)} igotzen nintzen,
\VUgras{bela-txalupakoa} \textsuperscript{(49)}, \VUgras{soken}
\textsuperscript{(52)} bidez \VUgras{mastara} \textsuperscript{(51)} igotzen
\VUgras{berga} \textsuperscript{(50)} arte, eta gero \VUgras{txalupatxoan}
\textsuperscript{(54)} jarraitzen nuen ibilaldia, \VUgras{arrantza-ontzi}
astunen artean \textsuperscript{(55)}, non \VUgras{sareak}
\textsuperscript{(56)} zabalduta lehortzen baitira.

%%K t10-17-1 p
17. --- \VUgras{Kaian} \textsuperscript{(58)} nire aurretik \VUgras{salgai}
mota guztiak \textsuperscript{(59)} igarotzen ziren, \VUgras{kutxetan}
\textsuperscript{(60)} edo \VUgras{fardeletan} \textsuperscript{(61)}
paketatuak. Gabarren \VUgras{zama} \textsuperscript{(63)} osatzen zuten, eta
\VUgras{garabi} indartsuek \textsuperscript{(62)} ateratzen zituzten eta
\VUgras{kamioietan} \textsuperscript{(64)} jartzen, \VUgras{garraiolariek}
haiek aitortu eta zerga \VUgras{aduanan} \textsuperscript{(65)} ordaintzen zuten
bitartean.

%%K t10-18-1 p
18. --- Azkenik, \VUgras{zubi altxagarrira} \textsuperscript{(66)} edo
\VUgras{esklusara} \textsuperscript{(69)} iritsita, \VUgras{dragaren} \textsuperscript{(68)}
ondotik igarota, zeinak \VUgras{ubidea} \textsuperscript{(67)} garbitzen
baitu, kai-muturrean atrakatzen nuen \VUgras{seinale-mastaren}
\textsuperscript{(70)} ondoan.

%%K t10-19-1 p
19. --- Kai-mutur horren beste aldetik atrakatzen dute hain zuzen
\VUgras{txalupek} \textsuperscript{(71)}, eskuadrako ofizialak lehorrera
ekartzen dituztenek. Ederra da ikustea nola altxatzen dituzten marinelek arraunak
batera, eta txalupa, \VUgras{olatuak} \textsuperscript{(72)} eramana, arinki
hurbiltzen den eskaileretara. Hemendik ikusten dira ondoen \VUgras{itsasgora} eta
\VUgras{itsasbeheraren} eta \VUgras{olatu-hausturaren} mugimendua
\VUgras{hondartzan} \textsuperscript{(73)}.

%%K t10-tit-7 sub
\VUsaut{3.0mm}
\VUpk{10.2pt}{40}{Ofizialen zirkulua.}
\VUsaut{3.0mm}

%%K t10-20-1 p
20. --- Etxera itzultzeko, \VUgras{tranbia elektrikoan} \textsuperscript{(74)}
igotzen nintzen, zeinaren \VUgras{geldialdia} \textsuperscript{(75)} ofizialen
zirkuluaren aurreko \VUgras{zutoinaren} ondoan baitago.

%%K t10-20-2 p
Zirkuluko \VUgras{balkoitik} \textsuperscript{(76)}, \VUgras{aingurez}
\textsuperscript{(77)}, kanoiez eta balaz apaindua, sarritan ikusten nuen
itsas-ofizialen gizarte distiratsua: \VUgras{tenienteak}
\textsuperscript{(78)} beren ezpatekin, \VUgras{artilleriako}
\textsuperscript{(79)} eta \VUgras{infanteriako} \textsuperscript{(80)}
\VUgras{kapitainak} beren \VUgras{sableekin}. Haien atzean \VUgras{marinel} bat
\textsuperscript{(81)} zegoen, zeinak \VUgras{teleskopioa} ematen baitzien
urrunean zer gertatzen zen ikusi nahi zutenean.

%%K t10-21-1 p
21. --- \VUgras{Guardiamarina} gazteak ere ikusi nituen
\textsuperscript{(82)}, beren \VUgras{komandantearengandik}
\textsuperscript{(83)} edo \VUgras{almirantearengandik} \textsuperscript{(84)}
aginduak jasotzen, \VUgras{kuartelmaisua} \textsuperscript{(85)} haiek ontzira
eramateko zain zegoen bitartean.

%%K t10-22-1 p
22. --- Balkoi horretatik ikuspegia bikaina da: hondartza osoa menderatzen du
bere \VUgras{denda} \textsuperscript{(86)} eta \VUgras{bainu-etxolekin}
\textsuperscript{(87)}, eta bainu-orduan ikusten da nola \VUgras{bainulariak}
\textsuperscript{(88)} alai jostatzen diren \VUgras{kasinoaren}
\textsuperscript{(89)} aurrean.

%%K t10-23-1 p
23. --- Hondartzaren amaieran itsasertzak \VUgras{penintsula} harritsu txiki bat
osatzen du \textsuperscript{(91)}, zeinaren kontra \VUgras{itsasbitsak}
\textsuperscript{(92)} hausten baitira eta zeinaren gainean \VUgras{itsasargia}
\textsuperscript{(93)} eraiki baitute, gauez marinelei bidea erakusten diena.
Penintsula hori lurrarekin oso \VUgras{lepo} estu batez baino ez dago lotuta
\textsuperscript{(90)}.

%%K t10-tit-8 sub
\VUsaut{3.0mm}
\VUpk{10.2pt}{40}{Itsasertzaren ikuspegi orokorra.}
\VUsaut{3.0mm}

%%K t10-24-1 p
24. --- \VUgras{Itsaslabarra} \textsuperscript{(94)} urrunago luzatzen da,
itsasoak ebakia, zeinak \VUgras{badien} \textsuperscript{(95)} eta
\VUgras{lurmuturren} segida bitxi bat zizelkatzen baitu, oso pintoresko eginez.

%%K t10-24-2 p
\VUgras{Lautadan} \textsuperscript{(96)}, zeinak \VUgras{itsasartea}
\textsuperscript{(98)} menderatzen baitu, oso \VUgras{gotorleku}
garrantzitsu bat \textsuperscript{(97)} eta txalet batzuk daude.

%%K t10-25-1 p
25. --- Itsasarte horrek lehorretik bereizten du oso \VUgras{uharte}
arriskutsu bat \textsuperscript{(99)}, \VUgras{harkaitzez} \textsuperscript{(100)}
eta \VUgras{arrezifez} inguratua, non batzuetan txalupak eta ontzi handiak ere
galtzen baitira. Laguntza zenbat eta bizkorragoa izan eta \VUgras{salbamendu-txalupak}
zenbat eta azkarrago iritsi, \VUgras{itsas hondoratze} horiek ikaragarriak dira, eta
ekinokzioko \VUgras{ekaitzetan} ontzi batzuk galdu dira tripulazio eta zama
guztiarekin.

%%K t10-25-2 p
Auzo hori izan arren, marinelek gogotik bisitatzen dute \VUgras{portu} hau.
\VUsaut{6.0mm}
\VUfilet{22mm}
